% language=us

% \nopdfcompression

% we have local fonts here so no: runpath=texruns:manuals/ontarget

% Timestamp: Porcupine Tree closure/contuation blu-ray and Peter Gabriel i/o
% in loop mode (December 2023).

\startcomponent ontarget-authentic

\environment ontarget-style

\logo[HP]     {HP}
\logo[OCE]    {OCE}
\logo[SUMATRA]{Sumatra}

\startchapter[title={Bitmap fonts}]

\def\PKxsize{0}
\def\PKysize{0}

\startluacode
    local readers = fonts.handlers.tfm.readers

    function document.loadPKshapes(name)
        local f = readers.loadpk(name)
        if not f then
            print("missing file @ 1")
        end
        local function set(f,target,char)
            local g = f and f.glyphs[char]
            local b = g and readers.showpk(g) or ""
            potrace.setbitmap(target,b)
        end
        set(f,"authentic-@", string.byte("@"))
        set(f,"authentic-i", string.byte("i"))
        set(f,"authentic-K", string.byte("K"))
        set(f,"authentic-æ", 26)
        set(f,"authentic-111", 111)
    end

    function document.showPKshapes(name)
        local f = readers.loadpk(name)
        if not f then
            print("missing file @ 2")
        end
        local g = f and f.glyphs[string.byte("m")]
        local b = g and readers.showpk(g) or ""
        tokens.setters.macro("PKxsize",g and g.xsize or 0)
        tokens.setters.macro("PKysize",g and g.ysize or 0)
        buffers.assign("bitmap",b)
    end
\stopluacode

\startMPdefinitions
    vardef showPKshape(expr _name, _polygon, _threshold) =
        save p ; path p ; p := lmt_potraced [
            threshold  = _threshold,
            stringname = _name,
            polygon    = _polygon,
            value      = "1",
            size       = 0,
        ] ;
        p shifted (0,-ypart llcorner p)
    enddef ;
\stopMPdefinitions

This is a follow up on potrace|-|generated outlines from bitmaps (see preceding
article, ``Unusual bitmaps''). Most of today's \TEX\ users have probably never
generated a document with bitmap fonts but back in the day we only had these,
almost always Computer Modern. Because we don't use bitmap fonts in \CONTEXT\
there is no need to support so|-|called \PK\ (bitmap) fonts in the backend, but,
because we did support them in \MKIV, we still have it in \LMTX. After all, what
is \TEX\ without bitmap fonts? We also owe it to Don Knuth.

It is a bit of challenge to load a traditional eight bit font using only a \TFM\
file because in \CONTEXT\ we map everything to \UNICODE. Regular \TYPEONE\ fonts
are still supported, although they are declared obsolete, and vendors have moved
on to \OPENTYPE\ fonts. When we define a font that has \TYPEONE\ resources, we
load it as if it were an \OPENTYPE\ font. These eight bit fonts thus become wide
(up to 64K glyphs) fonts internally. Normally we consult the \AFM\ and \PFB\
files and leave the \TFM\ files, if present at all, for what they are. You can
think of runtime \type {afmtotfm} conversion. An exception is the few traditional
math fonts that we support: Antykwa, Iwona and Kurier; here the \TFM\ files
provide dimensions.

For the purpose of demonstrating what comes next it is enough to know that in
principle one can still mess around with \TFM\ values, either outline or bitmap,
and that encoding files play a role in mapping them onto \UNICODE. We never set
out to do something like this, but, because we had the loaders available anyway,
some quick (few line) experiments of passing \PK\ bitmaps to the \METAPOST\
potrace helpers we got curious.

If you've run into an old \TEX\ document on the web you might have noticed bitmap
fonts being used. Often these are of a relatively low resolution. The reason that
one never noticed that in print is that, when read on screen, we see glyphs large
and can even zoom in, while in print they are seen small. When larger glyphs are
used (say in a title) the scaled glyphs are actually different bitmaps: they have
the same resolution as the smaller ones but more pixels because they are larger.
Take these characters at 600 dpi, scaled up from their original 10 point size:

\ctxlua{document.loadPKshapes("ontarget-cmr10.600pk")}

\starttexdefinition showPKshapes #1#2#3
    \startcombination[nx=4,ny=1,distance=2em]
        {\ruledhbox{\startMPcode fill showPKshape("authentic-@",#2,#3) scaled (2*600/#1); \stopMPcode}} {}
        {\ruledhbox{\startMPcode fill showPKshape("authentic-i",#2,#3) scaled (2*600/#1); \stopMPcode}} {}
        {\ruledhbox{\startMPcode fill showPKshape("authentic-K",#2,#3) scaled (2*600/#1); \stopMPcode}} {}
        {\ruledhbox{\startMPcode fill showPKshape("authentic-æ",#2,#3) scaled (2*600/#1); \stopMPcode}} {}
    \stopcombination
\stoptexdefinition

\startlocallinecorrection
    \texdefinition{showPKshapes}{600}{true}{1}
\stoplocallinecorrection

How do these patterns translate into an outline? For this we use the potrace
library that we have available in \LUAMETATEX\ and interfaced to \METAFUN\ in
\CONTEXT\ (as discussed in the companion article), although for fonts we follow a
more direct route.

\startlocallinecorrection
    \texdefinition{showPKshapes}{600}{false}{1}
\stoplocallinecorrection

\ctxlua{document.showPKshapes("ontarget-cmr10.600pk")}

So why do these glyphs look somewhat different (smoother) from a normal outline
Computer Modern? This is because the 600 may sound like a lot but actually isn't.
Distributed over an inch we have 600 pixels on 25.4 mm so roughly one pixel per
0.05 mm, which might be acceptable in a small print but not when scaled. Here is
how an \quote {m} in 600 dpi pixels is coded (\PKxsize\ by \PKysize\ pixels):

\start\switchtobodyfont[4pt]\typebuffer[bitmap]\stop

In the days of sending bitmaps to printers, possibly wrapped in a \POSTSCRIPT\
file, one could often recognize the characters from blobs like this. You might
also remember some of that line printer art. At any rate, it shows that 600 dots
per inch is much less than it sounds. A decent 2400 dpi (dots per inch) is more
normal these days for printing on presses (with ink) but although there are 1200
dpi laser printers, 600 became the norm. After all, toner particles are not that
small. For quite a while the authors used high speed \OCE\ low temperature toner
printer (first 508 dpi, later 600 dpi) and these were visually superior to most
of what the competition had \emdash\ but \OCE\ never quite managed to do the same
in color (only in lab testing, not reaching the market). Nowadays we use \HP\
full width high speed inkjet printers that give a rather good quality full color
experience at 600 dpi. This is what a scaled|-|up 2400 dpi bitmap look like:

\ctxlua{document.loadPKshapes("ontarget-cmr10.2400pk")}

\startlocallinecorrection
    \texdefinition{showPKshapes}{2400}{false}{1}
\stoplocallinecorrection

This already looks more crisp although we can go back to the smoother variant by
setting a higher threshold in potrace:

\startlocallinecorrection
    \texdefinition{showPKshapes}{2400}{false}{2}
\stoplocallinecorrection

We can go higher. In the next rendering we use 7200 dpi bitmaps and now we see
some details that didn't show up before. Keep in mind that subtle details might
not be noticed in 10 point running text. {\red Are the \quotation{ends of the at
and e cut?}}

\ctxlua{document.loadPKshapes("ontarget-cmr10.7200pk")}

\startlocallinecorrection
    \texdefinition{showPKshapes}{7200}{false}{1}
\stoplocallinecorrection

The top of the \quote {K} now has a little dent (likely not visible except with
magnification). If you ever viewed a document with Latin Modern outlines you
might recognize this. It is a side effect of outlines having that information
while a bitmap needs to get the exact bits, which won't happen if such a dent
stays beyond the pixel threshold. We now are ready for some real text. We define
two font features to load bitmap and outlines, respectively:

\startbuffer
\definefontfeature[whateverpk][default][reencode=ontarget-cmr.enc,bitmap=pk]
\definefontfeature[whateverpt][default][reencode=ontarget-cmr.enc,bitmap=outline]
\stopbuffer

\typebuffer \getbuffer

and some colors:

\startbuffer
\definecolor[pkcolor][r=1,t=.5,a=1]
\definecolor[ptcolor][b=1,t=.5,a=1]
\stopbuffer

\typebuffer \getbuffer

We use the following three font definitions in three overlaid examples (\in
{figure} [fig:pkdemo]). The results are close enough to justify a closer look at
the possibilities.

% plex has a large 20% more exheight

\startbuffer
\definefont[PKdemoA][file:lmroman10-regular.otf*default sa 1.2]
\definefont[PKdemoB][file:ontarget-cmr10.tfm*whateverpk sa 1.2]
\definefont[PKdemoC][file:ontarget-cmr10.tfm*whateverpt sa 1.2]
\stopbuffer

\typebuffer \getbuffer

\startplacefigure[title=Overlayed regular, bitmap and potraced samples.,reference=fig:pkdemo]
    \startoverlay
        {\ruledvbox{\hsize12cm\PKdemoA\relax  \samplefile{tufte}}}
        {\ruledvbox{\hsize12cm\PKdemoB\pkcolor\samplefile{tufte}}}
        {\ruledvbox{\hsize12cm\PKdemoC\ptcolor\samplefile{tufte}}}
    \stopoverlay
\stopplacefigure

In \in {figure} [fig:smooth] we see from top to bottom: Latin Modern OpenType
outlines, a bitmap Computer Modern and a potraced Computer Modern. The fourth
line has the bitmap and potraced overlaid. \in {Figure} [fig:smoothclip] shows a
small portion of the last one as seen on screen. Blown up, some drift is seen but
so far we didn't find a way to get rid of it. Some is due to the way glyph
streams get rendered and synchronized (after spacing, for instance).

\startplacefigure[title={Bitmap and potraced compared: \OPENTYPE, \PK, potraced & overlayed.},reference=fig:smooth]
    \scale[width=1tw]\bgroup
        \vbox\bgroup \forgetall\showglyphs \offinterlineskip
          \ruledhbox{\PKdemoA         Smooth}
          \vskip.1ex
          \ruledhbox{\PKdemoB\pkcolor Smooth}
          \vskip.1ex
          \ruledhbox{\PKdemoC\ptcolor Smooth}
          \vskip.1ex
          {\startoverlay
                {\ruledhbox{\pkcolor\PKdemoB Smooth}}
                {\ruledhbox{\ptcolor\PKdemoC Smooth}}
          \stopoverlay}
        \egroup
    \egroup
\stopplacefigure

\startplacefigure[title=An enlarged clip of the overlay.,reference=fig:smoothclip]
    \externalfigure[ontarget-authentic-001.png][height=.2th]
\stopplacefigure

When not overlaid we get this for a normal Latin Modern Regular:

\start\blank\PKdemoA\samplefile{tufte}\blank\stop

And this for a Computer Modern Roman \PK\ bitmap:

\start\blank\PKdemoB\samplefile{tufte}\blank\stop

The same Computer Modern Roman outlined by potrace gives this:

\start\blank\PKdemoC\samplefile{tufte}\blank\stop

Because these are outlines we can actually scale them nicely with
\type {\glyphscale 800} here:

\start\blank\PKdemoC \glyphscale 800 \samplefile{tufte}\blank\stop

We can additionally scale an extra \type {\glyphxscale 1200}. This can of course
also be done with bitmaps but outlines are a safer bet.

\start\blank\PKdemoC \glyphscale 800 \glyphxscale 1200 \samplefile{tufte}\blank\stop

The conversion happens in the backend and can have some impact on the runtime but
we cache the outlines, so a subsequent run is faster. Of course outlines are more
efficient than bitmaps in terms of bytes and viewers also tend to render them
better than bitmaps, which, especially at a lower resolution, can look pretty bad
(in some viewers).

One might wonder if bitmaps are worse than outlines. When we use a reasonable
resolution there is no need to generate more than one size, and in \CONTEXT\ we
assume this anyway because we scale all fonts to 10 big points and scale that
shared instance on demand. The 600 dpi \quote {m} that we showed has \PKxsize\
pixels in the horizontal direction. So, 75 of these (a line of text) need 4800
%\the \numexpr \PKxsize * 75 \relax\ pixels, which is more than enough for a
4\endash8 display. Unfortunately viewers still render a bitmap font somewhat
badly.

Let's dive a little into why bitmaps might render suboptimally in \PDF\ viewers.
Consider these three lines typeset in our three fonts:

\starttyping
\PKdemoA Smooth \vskip.1ex
\PKdemoB Smooth \vskip.1ex
\PKdemoC Smooth \vskip.1ex
\stoptyping

These \quote {Smooth} lines (this time in 12pt) end up in the \PDF\ file as
follows :

\starttyping
BT
/F1 10 Tf
1.195514 0 0 1.195514 0 30.316793 Tm [<000100020003>-28<000300040005>] TJ
/F2 10 Tf
1.195514 0 0 1.195514 0 15.415656 Tm [<010203>-28<030405>] TJ
/F3 10 Tf
1.195514 0 0 1.195514 0 0.514763 Tm [<010203>-28<030405>] TJ
ET
\stoptyping

This code first switches to font \type {/F1}, a wide outline font so we have
four|-|byte indices (in angle brackets). Next we trigger \type {/F2}, the bitmap
variant and finally the potraced \type {/F3}; these are both \TYPETHREE\ fonts so
they get two|-|byte indices. We don't scale except to the 10 big points design
size. After such a switch comes lines of text and there we do scale, here by
$1.195514$ in both directions. We're slightly off $1.2$ because the \PDF\ font
system (by tradition) is set up in PostScript (big) points so we need to scale up
a little from 12pt to 12bp. Scaling an outline is translated (in the end) to some
factor and the renderer can keep the device into account when it comes to
rounding. With bitmaps it's different, because these are not
mathematically|-|defined fonts, but some image that gets scaled. This can
introduce the first inaccuracy. An inline bitmap in \PDF\ is given between \type
{ID} and \type {EI} operators, as demonstrated in the next two charproc entries
for the \TYPETHREE\ bitmap font (reformatted and abridged to save space):

\starttyping
21 0 obj
<< /Length 40708 >>
stream
556 0 55 -22 498 703 d1
q
442 0 0 725 55 -22 cm
BI
  /W 442 /H 725 /IM true /BPC 1 /D [1 0]
ID ...bytes...
EI
Q
endstream
endobj
\stoptyping

Watch the different in length:

\starttyping
22 0 obj
<< /Length 42784 >>
stream
833 0 33 0 809 440 d1
q
775 0 0 440 33 0 cm
BI
  /W 775 /H 440 /IM true /BPC 1 /D [1 0]
ID ...bytes...
EI
Q
endstream
endobj
\stoptyping

In contrast, a character in the potraced outline font looks like this:

\starttyping
31 0 obj
<< /Length 3678 >>
stream
556 0 55 -22 498 703 d1
q
1 0 0 1 55 -22 cm
172 723.960456969 m 144.844454411 720.63823631 ...
Q
endstream
endobj
\stoptyping

The length is much smaller and the outline shape is just a sequence of moveto
(\type {m}), lineto (\type {l}) and curveto (\type {c}) operators mixed with
numbers.

\starttyping
32 0 obj
<< /Length 4260 >>
stream
833 0 33 0 809 440 d1
q
1 0 0 1 33 0 cm
68.5 434.441743787 m 32.2 431.509475939 1.9375 ...
Q
endstream
endobj
\stoptyping

Experiments demonstrated that it's better to use rounded widths because otherwise
(at least in \SUMATRA) we get some accumulated drift. The bitmap variants have a
transform matrix like this:

\starttyping
442 0 0 725 55 -22 cm
775 0 0 440 33   0 cm
\stoptyping

and the potraced outlines have:

\starttyping
1 0 0 1 55 -22 cm
1 0 0 1 33   0 cm
\stoptyping

This means that a bitmap again gets scaled, luckily by an integer, but still
there is some inaccuracy. In the end, we get the bits put on screen and
especially at small scales we end up with artifacts in positioning and scaling.
Where the font renderer is optimized for (indeed) rendering fonts, the bitmap
renderer isn't. In \in {figure} [fig:zoom] we see bitmaps being rendered bolder
when they become smaller, because in the end, even at high resolutions, we're not
talking pixels but bits (that can occupy multiple pixels). Outline fonts talk
pixels, bitmap fonts speak in bits.

\startplacefigure[title=Three zoom levels compared.,reference=fig:zoom]
    \startcombination[nx=3,ny=1]
        {\externalfigure[ontarget-authentic-002.png][height=3cm]} {smallest}
        {\externalfigure[ontarget-authentic-003.png][height=3cm]} {small}
        {\externalfigure[ontarget-authentic-004.png][height=3cm]} {normal}
    \stopcombination
\stopplacefigure

We haven't yet discussed how we got the bitmaps that we used. You might have
noticed in the examples that there are some differences with the outline when it
comes to dimensions. This is partly due to the fact that where Latin Modern is an
\OPENTYPE\ font with no limits to dimensions, Computer Modern has to accommodate
the limited number of heights, depths and widths that the \TFM\ format permits.
Think of arbitrary values of height compared to categories of height.

When you generate a bitmap you rely on scripts that do the work and these work
together with so|-|called printer modes as defined in the \METAFONT\ file \type
{modes.mf} (\hyphenatedurl {https://ctan.org/pkg/modes}). These modes are for
printers which means that there can be compensation going on: rounding up or down
of points exceeding bounding box edges, the size of printer pixels (toner, ink)
and accuracy of positioning them, etc. In our case, when we went for 8000 dpi we
ended up with a device that did more compensation than needed. Better is to
investigate time in figuring out how to control the machinery to cook up (maybe)
7200 dpi bitmaps because down the inclusion route there is some division by 72.
If we decide to play a bit more, we might as well first figure out how to control
the bitmap generation and see if we can come up with this 7200 resolution. Not
only does it divide nicely by 72 (for display) but also by 600 (for the average
printer). To what extent that matters is to be seen.

A bitmap font (instance) is generated by \METAFONT\ and driven by a (printer)
mode defined in \type {modes.mf}. We added this one:

\starttyping
mode_def potrace =
    mode_param (pixels_per_inch, 2 * 3600) ;
    mode_param (blacker, 0) ;
    mode_param (fillin, 0) ;
    mode_param (o_correction, 1) ;
    mode_common_setup_ ;
enddef;
\stoptyping

The \type {2 * 3600} is a trick to get around \METAFONT\ maxing out at 4096 but
internally being capable to deal with larger numbers. Of course we forgot to
run \typ {fmtutil-sys --byfmt mf} which is needed to get these modes in the
format file, but eventually we managed to generate the 7200 dpi \PK\ file for
\type {cmr10}. Generating is easiest done with \PDFTEX\ with \typ {\pdfmapfile
{}}, which wipes the mapping to a \TYPEONE\ file.

We mentioned using \METAPOST\ in our first attempts to get an idea how well
potrace can vectorize the bitmaps. Here is how that is done:

\starttyping
\startluacode
    local f = fonts.handlers.tfm.readers.loadpk("cmr10.pk")
    if f then
        local g = f.glyphs[string.byte("R")]
        if g then
            local b = fonts.handlers.tfm.readers.showpk(g)
            potrace.setbitmap("demo",b)
        end
    end
\stopluacode

\startMPpage[offset=1ts]
    draw lmt_potraced [
        stringname = name,
        value      = "1",
    ];
\stopMPpage
\stoptyping

In \in {figure} [fig:metapotrace] we demonstrate three such renderings. Watch how
the number of points increases as the shape gets better. In \in {figure}
[fig:metapotracecompare] we show the potraced variant alongside the \OPENTYPE\
outline. Left is the default potrace output, next comes the regular \OPENTYPE\
(here \CFF) outline, and at the right we see two potraced results, both with \typ
{optimize = true} passed; the first has the default tolerance of 0.2 and the
second uses 0.5 and thereby has less points. For sure the middle one has less
points which is nice, but the potrace ones are not that excessive so we can live
with it. We've run into cases (especially math) where the regular outlines are
also not perfect. Most will go unnoticed anyway given the small size at which
glyphs are normally rendered. When you look closely at the rightmost output
you'll notice that the bend in the stems makes for more points which is an
indication that the \METAFONT\ output might actually be more subtle. In \in
{figure} [fig:comparecontrols] we also show the controls and you see subtle
differences in the angles there. Also note that when there are no lines that
indicates that there is likely a lineto instead of a curveto.

\startluacode
    local function Example(filename,name,char)
        local f = fonts.handlers.tfm.readers.loadpk(filename)
        if not f then
            print("missing file @ 1")
        end
        local c = type(char) == "string" and string.byte(char) or tonumber(char)
        local g = f and c and f.glyphs[c]
        local b = g and fonts.handlers.tfm.readers.showpk(g) or ""
        potrace.setbitmap(name,b)
    end
    Example("ontarget-cmr10.600pk","demo-R-600","R")
    Example("ontarget-cmr10.2400pk","demo-R-2400","R")
    Example("ontarget-cmr10.7200pk","demo-R-7200","R")
    Example("ontarget-frualt.7200pk","demo-allrunes-7200",111)
\stopluacode

\startMPdefinitions
    vardef Example(expr name, clr, opt, tol, linestoo) =
        image (
            save p ; path p ; p := lmt_potraced [
                stringname = name,
                value      = "1",
                optimize   = opt,
                tolerance  = tol,
            ];
            p := p xsized 4cm ;
            draw p ;
            drawpoints       p withpen pencircle scaled 2.0 withcolor clr ;
        if linestoo :
            drawcontrollines p withpen pencircle scaled 0.5 withcolor clr ;
        fi
        )
    enddef ;
\stopMPdefinitions

\startplacefigure[title=Using \METAPOST\ for analysis.,reference=fig:metapotrace]
    \startMPcode
        draw Example("demo-R-600", "red",false,0.2,false) ;
        draw Example("demo-R-2400","red",false,0.2,false) shifted ( 5cm,0) ;
        draw Example("demo-R-7200","red",false,0.2,false) shifted (10cm,0) ;
        currentpicture := currentpicture xsized TextWidth ;
    \stopMPcode
\stopplacefigure

\startplacefigure[title=Comparing potraced bitmaps with \TYPEONE.,reference=fig:metapotracecompare]
\startMPcode
    picture p[] ;
    p[1] := Example("demo-R-7200","red",false,0.2,false) ;
    p[2] := lmt_outline [ text = "\definedfont[file:lmroman10-regular.otf]R", kind = "path" ] ;
    p[3] := Example("demo-R-7200","red",true,0.2,false) ;
    p[4] := Example("demo-R-7200","red",true,0.5,false) ;
    for i=1 upto 4 :
        p[i] := p[i] ysized 4cm ;
        p[i] := p[i] shifted ((i-1)*5cm,0) ;
    endfor
    draw p[1] ;
    draw p[3] ;
    draw p[4] ;
    draw p[2] withpen pencircle scaled 1/2 ;
    for i within p[2] :
        drawpoints       pathpart i withpen pencircle scaled 2.0 withcolor "blue" ;
    endfor
    currentpicture := currentpicture xsized TextWidth ;
\stopMPcode
\stopplacefigure

\startplacefigure[title=Comparing control points.,reference=fig:comparecontrols]
\startMPcode
    picture p[] ;
    p[1] := Example("demo-R-7200","red",true,0.5,true) ;
    p[2] := lmt_outline [ text = "\definedfont[file:lmroman10-regular.otf]R", kind = "path" ] ;
    p[1] := p[1] ysized 4cm ;
    p[2] := p[2] ysized 4cm ;
    p[1] := p[1] shifted (5cm,0) ;
    draw p[1] ;
    draw p[2] withpen pencircle scaled 1/2 ;
    for i within p[2] :
        drawpoints       pathpart i withpen pencircle scaled 2.0 withcolor "blue" ;
        drawcontrollines pathpart i withpen pencircle scaled 0.5 withcolor "blue" ;
    endfor
    currentpicture := currentpicture xsized TextWidth ;
\stopMPcode
\stopplacefigure

If you like the look of these shapes, take a look at Volume~E of Don Knuth's
\quotation {Computers \& Typesetting} series. An incredible amount of work went
into the details of the fonts and you'll run into brackets, beaks, crisps,
notches, slabs, juts, dishes and more elements and parameters that are used. For
instance the serifs as seen in the \quote {R} are actually made programmatically
(so that they can be discarded in the sans shapes). If you check out the proof
sheet of the \quote {R} you will only see the basic points that describe the
character, not the points we see above, after conversion to (any kind of)
outline.

So now that we can turn a \PK\ into a proper outline we're done, right? Well, not
entirely. Because we have relatively simple shapes (moveto, lineto and curveto)
we can directly go to \CFF\ and avoid the \METAPOST\ to \TYPETHREE\ conversions.
Because we already can load (and adapt) \CFF\ outlines it is not that complicated
to do the reverse. The backend already can include them so we can also borrow
code there.

When going from potrace output to \CFF, we need a high resolution, so we started
out again with 7200 dpi. Although in the end we were quite satisfied with the
results, we tested with 20000 dpi and thought it looks even better than the Latin
Modern successor (although one can argue about it). Some first experiments showed
that it was doable but it actually took a whole day to (test and) decide how to
get better results. For instance, \METAPOST\ uses floats while a renderer uses
integers, so we run into rounding issues if we delegate that (although we can
include floats in \CFF, it is not a success). When overlaying the \METAPOST\
output and the \CFF\ shapes the latter was quite disappointing: weird protrusions
and bubbles. Of course one may doubt the implementation, but double and triple
checking showed that we use the same numbers.

The results improved a lot when the results from potrace were multiplied by 10
(or 20) effectively giving very huge glyphs (on a 10000 unit canvas) and in the
page stream scaling down by that factor. By then using rounded values we got
enough precision to get the \CFF\ results close to the (always) high quality
\METAPOST\ rendering. In \in {figure} [fig:mpovercff] we can see how close they
became.

\startplacefigure[title={A \METAPOST\ rendering on top of the \CFF\ counterpart},reference=fig:mpovercff]
    \definefontfeature[MyFontA][default][reencode=ontarget-cmr.enc,bitmap=cff,resolution=7200]%
    \definedfont[file:ontarget-cmr10.tfm*MyFontA]%
    \scale[width=.3\textheight]{\ruledhbox{\startoverlay
      {\scale[height=5cm]{\hbox{\blue s}}}
      {\hbox
         xoffset -0.14mm
         yoffset  0.20mm
         \bgroup
           \startMPcode
             draw image (
               save p ; path p ; p := lmt_potraced [
                 filename   = "ontarget-cmr10.pk",
                 resolution = 7200,
                 index      = 115,
                 value      = "1",
               ] ;
               p := p ysized 5cm ;
               draw p ;
               drawpoints       p withpen pencircle scaled 2.00 withcolor "red" ;
               drawcontrollines p withpen pencircle scaled 0.25 withcolor "red" ;
             ) scaled 1.06 ;
           \stopMPcode
       \egroup}
    \stopoverlay}}
\stopplacefigure

The next question is, how do we control this: \PK, \METAPOST\ \TYPETHREE\ or
native \CFF ? Even more challenging is that we had wanted to use different
resolutions, and mix these three methods, if only because we want to be able to
experiment and document the mix. That means that it has to be available not only
in the font feature mechanism, which is rather trivial, but also in the backend,
which then involves a font with the same name (say \type {CMR10}) to be rendered
differently, so we need different handlers, distinctive caching of streams, etc.
In the end we got there. It is even possible to mix variants with different
potrace parameters in one document.

We start with \CFF\ definitions. In \in {figure} [fig:cffoverlayed] we show three
resolutions overlaid, with the 7200 dpi variant below. In \in {figure}
[fig:mpovercffagain] we show the default and optimized (fewer points) traces.

\startbuffer
\definefontfeature[CMRCFF][reencode=ontarget-cmr.enc,bitmap=cff]

\definefontfeature[MyFontCffA][default,CMRCFF][resolution=7200]
\definefontfeature[MyFontCffB][default,CMRCFF][resolution=2400]
\definefontfeature[MyFontCffC][default,CMRCFF][resolution=600]
\definefontfeature[MyFontCffD][default,CMRCFF][resolution=7200,potrace={optimize=true}]
\stopbuffer

\getbuffer \typebuffer

The overlays look fuzzy, demonstrating the need for high resolutions. Font
definitions are done as follows; we only show one definition:

\starttyping
\definefont[MyFont][file:ontarget-cmr10.tfm*MyFontCffA]
\stoptyping

\startplacefigure[title={The 7200, 2400 and 600 \CFF\ variants overlayed.},reference=fig:cffoverlayed]
    \startoverlay
        {\vbox{\colored[r=1,a=1,t=.5]{\definedfont[file:ontarget-cmr10.tfm*MyFontCffA]\samplefile{tufte}}}}
        {\vbox{\colored[g=1,a=1,t=.5]{\definedfont[file:ontarget-cmr10.tfm*MyFontCffB]\samplefile{tufte}}}}
        {\vbox{\colored[b=1,a=1,t=.5]{\definedfont[file:ontarget-cmr10.tfm*MyFontCffC]\samplefile{tufte}}}}
    \stopoverlay
    \blank
    \vbox{\definedfont[file:ontarget-cmr10.tfm*MyFontCffA]\samplefile{tufte}}
\stopplacefigure

\startplacefigure[title={The normal 7200 and optimized 7200 dpi variants},reference=fig:mpovercffagain]
    \startoverlay
        {\vbox{\colored[r=1,a=1,t=.5]{\definedfont[file:ontarget-cmr10.tfm*MyFontCffA]\samplefile{tufte}}}}
        {\vbox{\colored[b=1,a=1,t=.5]{\definedfont[file:ontarget-cmr10.tfm*MyFontCffD]\samplefile{tufte}}}}
    \stopoverlay
    \blank
    \vbox{\definedfont[file:ontarget-cmr10.tfm*MyFontCffA]\samplefile{tufte}}
\stopplacefigure

\startbuffer
\definefontfeature[CMRPK][reencode=ontarget-cmr.enc,bitmap=pk]

\definefontfeature[MyFontBitmapA][default,CMRPK][resolution=7200]
\definefontfeature[MyFontBitmapB][default,CMRPK][resolution=2400]
\definefontfeature[MyFontBitmapC][default,CMRPK][resolution=600]
\stopbuffer

\getbuffer \typebuffer

These definitions are used in \in {figure} [fig:pkoverlayed]. The differences
aren't immediately apparent in small print but zoom in (if you're online) and
you'll understand why we need more than 600 dpi to feel comfortable.

\startplacefigure[title={The 7200, 2400 and 600 \PK\ variants overlayed.},reference=fig:pkoverlayed]
    \startoverlay
        {\vbox{\colored[r=1,a=1,t=.5]{\definedfont[file:ontarget-cmr10.tfm*MyFontBitmapA]\samplefile{tufte}}}}
        {\vbox{\colored[r=1,a=1,t=.5]{\definedfont[file:ontarget-cmr10.tfm*MyFontBitmapB]\samplefile{tufte}}}}
        {\vbox{\colored[b=1,a=1,t=.5]{\definedfont[file:ontarget-cmr10.tfm*MyFontBitmapC]\samplefile{tufte}}}}
    \stopoverlay
    \blank
    \vbox{\definedfont[file:ontarget-cmr10.tfm*MyFontBitmapA]\samplefile{tufte}}
\stopplacefigure

\startbuffer
\definefontfeature[CMROUTLINE][reencode=ontarget-cmr.enc,bitmap=outline]

\definefontfeature[MyFontOutlineA][default,CMROUTLINE][resolution=7200]
\definefontfeature[MyFontOutlineB][default,CMROUTLINE][resolution=2400]
\definefontfeature[MyFontOutlineC][default,CMROUTLINE][resolution=600]
\stopbuffer

\getbuffer \typebuffer

Finally we show the \METAPOST|-|generated threesome in \in {figure}
[fig:outlineoverlayed] and these look quite reasonable. One thing to keep in mind when
wrapping shapes into a \TYPETHREE\ font is that one has to make sure that color keeps
working.

\startplacefigure[title={The 7200, 2400 and 600 \METAPOST\ variants overlayed.},reference=fig:outlineoverlayed]
    \startoverlay
        {\vbox{\colored[r=1,a=1,t=.5]{\definedfont[file:ontarget-cmr10.tfm*MyFontOutlineA]\samplefile{tufte}}}}
        {\vbox{\colored[r=1,a=1,t=.5]{\definedfont[file:ontarget-cmr10.tfm*MyFontOutlineB]\samplefile{tufte}}}}
        {\vbox{\colored[b=1,a=1,t=.5]{\definedfont[file:ontarget-cmr10.tfm*MyFontOutlineC]\samplefile{tufte}}}}
    \stopoverlay
    \blank
    \vbox{\definedfont[file:ontarget-cmr10.tfm*MyFontOutlineA]\samplefile{tufte}}
\stopplacefigure

So, how useful is all this? Maybe it's time for a revival of \METAFONT. Or maybe
we can get some old designs out of the archives where they got tagged obsolete
and use them again. Or maybe it's just for the fun of it. We started out with
\PK\ bitmaps that we need to support anyway. Next we were curious how well a
traced outline would look, and for that using a \METAPOST\ \TYPETHREE\ font makes
sense. Then we took the challenge to turn potrace output into a \CFF\ \OPENTYPE\
font. We could now make a whole tool chain but it makes little sense: we can do
it in \CONTEXT, so we're fine, and we don't expect widespread usage outside the
\TEX\ ecosystem. Next on the agenda is to locate some interesting
\hbox{\METAFONT}s and see if they can be put to use. In case you wonder how these
eight bit fonts fit into the \UNICODE\ ecosystem, don't worry, we can use the
virtual font mechanism to hook shapes into existing fonts (which is what we do
anyway) or create combined fonts. We can even make variable fonts using
\METAFONT, but for that we need to become experienced designers first. As a
teaser we show a character from the runes font by Carl|-|Gustav Werner in \in
{figures} [fig:allrunes] \in {and} [fig:allrunesclip]. Mikael, knowing the
author, promised to come up with a proper encoding for that one, so stay tuned.

\definefontfeature[MyRunesA][reencode=ontarget-frualt.enc,pfbfile=ontarget-frualt.pfb]
\definefontfeature[MyRunesB][reencode=ontarget-frualt.enc,bitmap=cff,resolution=7200]
\definefontfeature[MyRunesC][reencode=ontarget-frualt.enc,bitmap=pk,resolution=7200]
\definefontfeature[MyRunesD][reencode=ontarget-frualt.enc,bitmap=outline,resolution=7200]

% \ctxlua{document.loadPKshapes("ontarget-frualt.pk")}
%
% actually there is no need to go through mp any more as we can mix them in the backend
%
% \startplacefigure[title={A simple test with one of the allrunes fonts. Left is using the \PK\ at a 7200 resolution. Next is the potraced 7200 \PK\ original outline. Third from left is the generated \CFF. Right is the shipped \PFB.},reference=fig:allrunes]
%     \startcombination[nx=4,ny=1,distance=2em]
%       {\ruledhbox{\startMPcode
% `      % save p ; picture p; p := image(fill showPKshape("authentic-111",true,1)) ;
%          save p ; picture p; p := textext("\setbox\scratchbox\hbox{\definedfont[file:ontarget-frualt.tfm*MyRunesC]a}\dp\scratchbox\zeropoint\box\scratchbox") ;
%          p := p shifted -llcorner p ;
%          draw p ysized 4cm ;
%       \stopMPcode}} {\TYPETHREE: \PK}
%       {\ruledhbox{\startMPcode
%        % save p ; picture p; p := Example("demo-allrunes-7200","darkred",false,0.2,false) ;
%          save p ; picture p; p := textext("\setbox\scratchbox\hbox{\definedfont[file:ontarget-frualt.tfm*MyRunesD]a}\dp\scratchbox\zeropoint\box\scratchbox") ;
%          p := p shifted -llcorner p ;
%          draw p ysized 4cm ;
%       \stopMPcode}} {\TYPETHREE: \METAPOST}
%       {\ruledhbox{\startMPcode
%          save p ; picture p; p := textext("\setbox\scratchbox\hbox{\definedfont[file:ontarget-frualt.tfm*MyRunesB]a}\dp\scratchbox\zeropoint\box\scratchbox") ;
%          p := p shifted -llcorner p ;
%          draw p ysized 4cm ;
%       \stopMPcode}} {\OPENTYPE: \CFF}
%       {\ruledhbox{\startMPcode
%          save p ; picture p; p := textext("\setbox\scratchbox\hbox{\definedfont[file:ontarget-frualt.tfm*MyRunesA]a}\dp\scratchbox\zeropoint\box\scratchbox") ;
%          p := p shifted -llcorner p ;
%          draw p ysized 4cm ;
%       \stopMPcode}} {\TYPEONE: \PFB}
%     \stopcombination
% \stopplacefigure

\protected\def\MySample#1#2#3%
  {\scale[height=#1cm]\bgroup
     \setbox\scratchbox\hbox{\definedfont[file:ontarget-frualt.tfm*MyRunes#2]a}%
     \dp\scratchbox\zeropoint
     \ht\scratchbox\dimexpr\ht\scratchbox+#3pt\relax
     \box\scratchbox
   \egroup}

\startplacefigure[title={A simple test with one of the allrunes fonts. Left is using the \PK\ at a 7200 resolution. Next is the potraced 7200 \PK\ original outline. Third from left is the generated \CFF. Right is the shipped \PFB.},reference=fig:allrunes]
    \startcombination[nx=4,ny=1,distance=2em]
        {\ruledhbox{\MySample{4}{C}{0}}} {\TYPETHREE: \PK}
        {\ruledhbox{\MySample{4}{D}{0}}} {\TYPETHREE: \METAPOST}
        {\ruledhbox{\MySample{4}{B}{0}}} {\OPENTYPE : \CFF}
        {\ruledhbox{\MySample{4}{A}{0}}} {\TYPEONE  : \PFB}
    \stopcombination
\stopplacefigure

\startplacefigure[title={Enlarged clips the four variants in \in {figure} [fig:allrunes].},reference=fig:allrunesclip]
    \startcombination[nx=4,ny=1,distance=1em]
        {\clip[ny=6,y=1,nx=5,x=2]{\MySample{20}{C}{.075}}} {\TYPETHREE: \PK}
        {\clip[ny=6,y=1,nx=5,x=2]{\MySample{20}{D}{.075}}} {\TYPETHREE: \METAPOST}
        {\clip[ny=6,y=1,nx=5,x=2]{\MySample{20}{B}{.075}}} {\OPENTYPE : \CFF}
        {\clip[ny=6,y=1,nx=5,x=2]{\MySample{20}{A}{.075}}} {\TYPEONE  : \PFB}
    \stopcombination
\stopplacefigure

\stopchapter

\stopcomponent
